\chapter{Monitoring an attack while it runs}\label{ch:monitoring}

An analyst who simulates an attack on a release can look at the results whenever
she likes. If the bound she computes at each look is exact, the argument runs, then
whatever she concludes at that look is exact as well. Looking more often uses more
of the evidence and never less of it. If an audit fails narrowly, a second audit on
freshly drawn targets is new and independent evidence, and it deserves its own
hearing. Each step of this argument holds up when it is examined alone. The
one-sided Clopper--Pearson bound of Chapter~\ref{ch:estimate} has exactly the
coverage it claims at the sample size where it is computed, and a fresh draw of
targets owes nothing to the draw before it. An analyst who worked this way would
be doing what careful people do with any measurement that comes in slowly, and a
committee that read her record would find a valid bound beside the decision to
release.

The argument fails, and the size of the failure is easy to measure. Take a release
whose true risk is 0.201 against a tolerance of 0.2, so that any certificate it
receives is false. % source: RESEARCH_BRIEF.md; scripts/injector_experiment.py (audit_restarts)
Audit it with 500 simulated attacks and a one-sided bound at the 5 per cent level,
and the bound passes in about 3.5 per cent of audits.
% source: artifacts-injector/results.json, audit_restarts, attempts=1, exact_rate 0.03497
Now let the analyst try again on fresh targets after each failure, up to 20
attempts, each judged as a new test at the 5 per cent level. Over 20,000 simulated
programmes the release was certified in 51.0 per cent of them. Had the same 5 per
cent been divided among the 20 attempts before the first one ran, the rate would
have been 3.9 per cent.
% source: artifacts-injector/results.json, audit_restarts, attempts=20: reset_alpha rate 0.50975 (95% interval 0.50280 to 0.51670), preallocated_budget rate 0.03875 (0.03612 to 0.04152); RESEARCH_BRIEF.md
Every single bound in the first procedure was exact. The decision built from them
was a coin toss.

At the Bank, the problem arrived in the ordinary way. Naledi had been asked to
audit the full pooled attack on the credit register at a minimum cell count of 150,
the attack in which the adversary holds all six monthly vintages, a reading of the
loan's exposure in each, and the loan's history. The chair wanted the result for
the committee on Thursday. Naledi started the simulation on Monday evening, and on
Tuesday morning she opened the output to see how it was going. Nothing in her
training told her that this was a decision.

This chapter answers one question. What must an audit fix before it starts, and
what may it leave open, so that the analyst can look at the evidence as it arrives,
stop when it is clear and try again when it is not, and still hand the committee a
certificate whose error probability is the one written on it?

\section{Why valid looks add up to an invalid decision}\label{c06:sec-retry}

The retry experiment has a short explanation. Write $A_j$ for the event that
attempt $j$ certifies a release whose true risk is at or above the tolerance. Each
attempt, taken alone, has $\Prob(A_j) \le \alpha$. The committee, however, does not
act on one attempt. It acts on the event that some attempt certified, which is the
union $A_1 \cup \dots \cup A_n$, and a union of events can be far more probable
than any one of them.

\begin{proposition}[Retrying an audit]\label{c06:prop-retry}
Let a release have true risk $r \ge \tau$, and let an analyst run up to $n$ audits,
stopping at the first that certifies.
\begin{enumerate}
\item If the audits use independent draws of targets and each certifies with
probability $f > 0$, the probability that the release is certified is
$1 - (1-f)^n$, which tends to one as $n$ grows.
\item If audit $j$ is run at a level $\alpha_j$ fixed before the first audit begins,
with $\sum_{j=1}^n \alpha_j \le \alpha$, the probability that the release is
certified is at most $\alpha$, whatever the dependence between the audits and
whatever rule the analyst uses to decide whether to run another.
\end{enumerate}
\end{proposition}

The first part reproduces the repository's figures. With 500 attacks and a
tolerance of 0.2, the 95 per cent bound passes only when the count of successes is
small enough, and at a true risk of 0.201 that happens with probability $f \approx
0.035$. The single-audit rate sits below 5 per cent because the binomial count is
discrete and the exact bound is conservative. Twenty independent attempts give
$1 - (1-0.035)^{20} \approx 0.51$, which is what the simulation found.
% source: artifacts-injector/results.json: exact_rate 0.03497 (one attempt), 0.50934 (20 attempts, reset alpha)
The second part is the remedy, and it rests on nothing more than the union bound.
Its condition is that the levels are written down before any outcome is seen. A
level chosen after the third failure, once the analyst has a feeling about the
fourth, is a level chosen with the data, and the union bound no longer applies to
it.

Figure~\ref{c06:fig-retry} shows the two rules side by side at one, five and twenty
attempts. With one attempt they coincide. With five, resetting the level each time
already certifies an unsafe release in 16.5 per cent of programmes, against 4.1
per cent when the level is divided in advance. With twenty the gap is the one
reported above. The pre-allocated rule stays close to its single-attempt rate
however many attempts it allows, because the extra attempts are paid for out of a
budget that does not grow.

\begin{figure}[ht]
\centering
\begin{tikzpicture}
\begin{axis}[
  ybar, bar width=14pt, width=0.78\textwidth, height=6cm,
  symbolic x coords={1,5,20}, xtick=data,
  xlabel={Audit attempts allowed}, ylabel={False certification (per cent)},
  ymin=0, ymax=70, enlarge x limits=0.25,
  legend style={at={(0.03,0.97)},anchor=north west,draw=none,font=\small},
  area legend, nodes near coords, nodes near coords style={font=\scriptsize},
  every node near coord/.append style={/pgf/number format/fixed,/pgf/number format/precision=1}]
\addplot[fill=gray!70, draw=black] coordinates {(1,3.6) (5,16.5) (20,51.0)};
\addplot[fill=white, draw=black, postaction={pattern=north east lines}] coordinates {(1,3.6) (5,4.1) (20,3.9)};
\legend{Level reset to 5\% at each attempt, 5\% divided among attempts in advance}
\end{axis}
\end{tikzpicture}
% source: artifacts-injector/results.json, audit_restarts (rates 0.03595, 0.16525, 0.50975; 0.03595, 0.0409, 0.03875)
\caption{Retrying a failed audit. A release with true risk 0.201 is audited against
a tolerance of 0.2 with 500 attacks per attempt, so every certificate is false.
Solid dark bars reset the level to 5 per cent at each attempt; hatched light bars
divide 5 per cent among the attempts before the first one runs. Each bar is a share
of 20,000 simulated programmes.}
\label{c06:fig-retry}
\end{figure}

Two features of the experiment deserve a plain statement. First, the analyst in it
is honest. She does not fabricate an outcome or hide a bound; she runs a valid audit,
and when it fails she runs another valid audit. The error is in the record, which
shows the last audit and not the attempts before it. Second, the experiment
assumes independent attempts, which is the most favourable case for the analyst.
If the attempts share anything, such as a pool of targets from which the failures
are quietly removed, the inflation can be worse, and Section~\ref{c06:sec-process}
returns to that.

The same mechanism works inside a single run. An analyst who computes the
fixed-sample bound after every batch of attacks and stops at the first batch where
it passes has also turned one test into many, though the tests are now dependent,
because each look contains the looks before it. The repository measures this on the
credit register. The reference attack is the pooled attack against a release that
keeps at least five loans in every cell, whose success rate is 0.269.
% source: paper/generated/numbers.tex, \runseqrisk 0.269
An institution simulates 4,000 attacks and computes the 95 per cent bound 20 times
along the way. In 500 such runs, the bound fell below the true risk at some look in
21 per cent of them. Any of those runs could have ended in a release on a bound that
was wrong.
% source: paper/generated/numbers.tex, \runseqmonitorreplicates 4,000, \runseqlooks 20, \runseqruns 500, \runseqfixederror 21
The stated level of a fixed-sample bound is a statement about one sample size
chosen in advance. With unlimited looks the inflation does not stop at 21 per cent:
the law of the iterated logarithm implies that a fixed-sample bound, recomputed
forever, eventually falls below the truth with probability one. Online experiments
in commerce met the same problem when analysts watched their tests daily and
stopped on a good day \citep{johari2022}.

On Tuesday morning Naledi's bound for the floor of 150 stood just above the
tolerance. On Wednesday it stood just below. The chair asked whether Wednesday's
figure was the answer. Naledi could not say yes, because the 95 per cent printed
beside it referred to a number of attacks that nobody had fixed, and she had
already looked once.

\section{Betting against the hypothesis that the release is unsafe}\label{c06:sec-martingale}

The way out begins by asking the question the other way round. The committee wants
to release only when there is strong evidence that the risk lies below the
tolerance. The hypothesis to be rejected is therefore that the release is unsafe,
\[
H : r \ge \tau ,
\]
and a certificate is a rejection of $H$. The tool that allows looking at will is a
measure of evidence against $H$ that cannot grow large, on average, while $H$ is
true. A betting game gives one.

Let $X_b \in \{0,1\}$ record whether the $b$th simulated attack succeeds, each on a
freshly drawn target. Before attack $b$ the auditor chooses a number
$q_b \in [0, \tau)$, using only what she has seen so far. She starts with wealth
$E_0 = 1$ and updates it after each attack by
\begin{equation}\label{c06:eq-bet}
E_b = E_{b-1} \Big(\frac{q_b}{\tau}\Big)^{X_b}
\Big(\frac{1-q_b}{1-\tau}\Big)^{1-X_b} .
\end{equation}
The factor is the ratio of the probability of the observed outcome under a
success rate of $q_b$ to its probability under a success rate of $\tau$, as in the
sequential test of \citet{wald1945}, with an alternative that may change from draw
to draw. A failed attack multiplies the wealth by $(1-q_b)/(1-\tau) > 1$; a
successful one multiplies it by $q_b/\tau < 1$. The auditor is betting that attacks
fail more often than $H$ allows, and she gains when they do.

If $H$ is true, the bet is unfavourable to her. For a success probability $p \ge
\tau$ and any $q < \tau$, the expected factor is
\begin{equation}\label{c06:eq-factor}
p\,\frac{q}{\tau} + (1-p)\,\frac{1-q}{1-\tau}
= 1 + \frac{(p-\tau)(q-\tau)}{\tau(1-\tau)} \le 1 ,
\end{equation}
because $p - \tau \ge 0$ and $q - \tau < 0$. Her wealth then has the structure of
a gambler's fortune in a game that does not favour her.

\begin{definition}[Test supermartingale]\label{c06:def-test}
A sequence $(E_b)_{b \ge 0}$ of nonnegative random variables, adapted to the
information $(\mathcal{F}_b)$ available after each attack, is a test
supermartingale for $H$ if $E_0 = 1$ and
$\E[E_b \mid \mathcal{F}_{b-1}] \le E_{b-1}$ for every $b$ whenever $H$ is true.
\end{definition}

The process \eqref{c06:eq-bet} is a test supermartingale for $H$ provided each
attack, given everything recorded before it, succeeds with probability at least
$\tau$ when $H$ is true, and provided each $q_b$ is fixed before $X_b$ is seen.
The first condition concerns how targets are drawn, and it is not automatic;
Section~\ref{c06:sec-process} shows how it fails. The second is the rule that a bet
is placed before the card is turned. Under both, a classical inequality limits how
often a test supermartingale can ever become large.

\begin{proposition}[Ville's inequality]\label{c06:prop-ville}
If $(E_b)$ is a test supermartingale for $H$ and $H$ is true, then for every
$c > 0$,
\begin{equation}\label{c06:eq-ville}
\Prob\Big\{\sup_{b \ge 0} E_b \ge c\Big\} \le \frac{1}{c} .
\end{equation}
\end{proposition}

The inequality is due to \citet{ville1939}, and its force lies in the supremum.
The bound holds for the largest value the wealth ever reaches, so it holds at the
moment the auditor chooses to stop, whatever rule she used to choose it. The audit
rule follows at once: certify the release the first time $E_b \ge 1/\alpha$, and
never otherwise. The chance of certifying an unsafe release is then at most
$\alpha$, and the auditor may look after every attack, stop when the evidence is
clear, or continue when a reviewer asks for more, without changing that figure.
Looking costs nothing in validity, because the guarantee was never tied to a
sample size.

The simplest bet shows what the guarantee costs. With $q_b = 0$ throughout, each
failure multiplies the wealth by $1/(1-\tau)$ and the first success sets it to zero
for good. After $B$ failures $E_B = (1-\tau)^{-B}$, and the audit certifies once
$B \ge \log(1/\alpha)/\{-\log(1-\tau)\}$. That is the zero-success count of the
fixed-sample bound in Chapter~\ref{ch:estimate}. At a tolerance of 0.12 and
$\alpha = 0.05$ it is 24 attacks.
% source: computed with privacy_risk.adaptive_audit.zero_success_budget(0.12, 0.05, 1.0) = 24; formula paper/publication_versions/jpc_main.tex, Proposition prop:cost
For an audit that expects no successes, monitoring is free. The catch is in the
last clause of the rule: a single success ends the audit, and an attack that
succeeds even rarely will produce one. A bet that leaves room for successes has to
be placed at $q_b > 0$, and then each failure earns less.

\section{Confidence sequences}\label{c06:sec-cs}

A committee often wants a number as well as a verdict: an upper bound on the risk
that it can report, which remains valid however long the simulation ran and however
often someone looked at it. A confidence sequence supplies one. Instead of one
betting process for one tolerance, take one for every candidate risk $r$ and mix the
alternative over all values with a beta prior. With $S_B$ successes in $B$ attacks,
the mixed process is
\begin{equation}\label{c06:eq-mixture}
E_B(r) = \frac{\mathrm{B}(a + S_B,\ b + B - S_B)}
{\mathrm{B}(a,b)\, r^{S_B} (1-r)^{B - S_B}} ,
\end{equation}
where $\mathrm{B}$ is the beta function and $(a,b)$ are the prior's parameters
\citep{howard2021,ramdas2023}.

\begin{proposition}[Anytime-valid confidence sequence]\label{c06:prop-cs}
If the simulated attacks are conditionally Bernoulli with risk $r_0$, then
$E_B(r_0)$ is a nonnegative martingale with $E_0(r_0) = 1$, and the sets
\begin{equation}\label{c06:eq-cs}
\mathcal{C}_B = \{ r \in (0,1) : E_B(r) < 1/\alpha \}
\end{equation}
contain $r_0$ at every $B$ simultaneously with probability at least $1 - \alpha$.
\end{proposition}

The upper end of $\mathcal{C}_B$ is a bound the analyst may compute after every
attack. In the monitoring study above, the confidence sequence fell below the true
risk at some look in none of the 500 runs, against 21 per cent for the fixed-sample
bound.
% source: paper/generated/numbers.tex, \runseqanytimeerror 0, \runseqfixederror 21
The price is a wider interval. On the same counts the fixed-sample upper bound
averaged 0.281 and the anytime bound 0.296, so the anytime bound was larger by a
factor of about 1.05.
% source: paper/generated/numbers.tex, \runseqfixedupper 0.281, \runseqanytimeupper 0.296, \runseqwidthratio 1.05
For a release whose risk sits well away from its tolerance, that difference
changes nothing.

Near zero risk the price is larger, and Table~\ref{c06:tab-anytime} gives it. To
certify a tolerance of one in a hundred at the 5 per cent level with no observed
success, a sample size fixed in advance needs 299 attacks, and the confidence
sequence with a uniform prior needs 984. For one in a thousand the figures are 2,995
and 12,417. The ratio, between three and four, is not the price of monitoring.
The zero bet of the previous section is also valid under monitoring, and it
certifies after the fixed-sample count. What the larger count buys is the ability
to keep going after a success. The mixture spends part of its evidence on
alternatives at which successes are plausible, so a single success does not end the
audit. An institution that expects some successes, as it will whenever the risk is
not negligible, needs that ability and pays for it.

\begin{table}[ht]
\centering
\caption{Successless attacks needed to certify a tolerance at the 5 per cent level.
The fixed-sample count holds at one sample size chosen in advance; the anytime
count holds while the run is watched, using the confidence sequence
\eqref{c06:eq-cs} with a uniform prior.}
\label{c06:tab-anytime}
% source: paper/generated/table_anytime.tex; paper/generated/numbers.tex (\runseqcertfixed001, \runseqcertanytime001, \runseqcertfixed0001, \runseqcertanytime0001)
\begin{tabular}{rrrr}
\toprule
Tolerance $\tau$ & Fixed sample & Anytime & Ratio \\
\midrule
0.01 & 299 & 984 & 3.3 \\
0.001 & 2{,}995 & 12{,}417 & 4.1 \\
\bottomrule
\end{tabular}
\end{table}

The confidence sequence and the betting rule of the previous section answer
different questions, and a record should not mix them. The upper end of
$\mathcal{C}_B$ is a bound on the risk. The wealth $E_b$ of \eqref{c06:eq-bet} is
evidence against one hypothesis at one tolerance; neither it nor its reciprocal is
an estimate of risk, and reporting it as one is an error. A decision record can
carry both, in separate fields, each with its level.

Naledi switched the Thursday audit to the confidence sequence and started it again
from zero on fresh targets, recording that she had done so and why. She could now
look each morning without spending anything. What she could not do was erase
Tuesday, and the next section is about how a programme keeps account of Tuesdays.

\section{The registered sequential audit}\label{c06:sec-gate}

A single release against a single attack is the simplest case. A programme meets a
harder one. The list of attacks worth testing is not fixed when the work begins:
analysts read one audit and propose another, reviewers ask for a pooled version of
an attack already run, and Chapter~\ref{ch:programmes} shows why pooled attacks
need audits of their own. A procedure that controls error for a family fixed in
advance does not cover a family that grows in response to its own results. The
journal article that accompanies this book sets out an audit for that situation,
and this section gives its structure. The error control is the online form of the
Bonferroni rule \citep{foster2008,tian2021}, applied to betting processes of the
kind in \eqref{c06:eq-bet}; what is specific to disclosure is the object being
audited and the conditions on how its targets are drawn.

Each pair of a candidate release and a declared attack is a \emph{coordinate}. A
coordinate is registered before its first outcome, and registration fixes four
things: the release and attack, the tolerance $\tau_i$, the attack code and target
law, and a share $w_i \ge 0$ of the programme's error budget. The shares must
satisfy $\sum_i w_i \le 1$ on every path, counting coordinates that are later
withdrawn. Coordinate $i$ runs its own betting process $E_{i,b}$ and is certified the
first time $E_{i,b} \ge (\alpha w_i)^{-1}$, after which it accepts no more
outcomes. A release is eligible only when every coordinate declared for it has been
registered and certified; an attack that has not been audited is not treated as
having zero risk.

\begin{assumption}[Conditional attack validity]\label{c06:as-fresh}
For each registered coordinate $i$ whose true risk satisfies $r_i \ge \tau_i$, and
for each of its draws $b$, the conditional probability that the attack succeeds,
given everything recorded before the draw on every coordinate, is at least
$\tau_i$.
\end{assumption}

The assumption holds when each draw attacks a fresh target from the registered
target law, with fresh noise in the adversary's information, independently of
everything recorded so far. It is stronger than a statement about the average
success rate, and Section~\ref{c06:sec-process} shows designs that break it.

\begin{proposition}[Familywise control under adaptive registration]\label{c06:prop-gate}
Suppose each coordinate's release, attack, tolerance and share are fixed at its
registration, with $w_i = 0$ for a coordinate never registered; the coordinate to
be audited at each step is chosen from what has already been recorded; each bet
$q_{i,b} \in [0, \tau_i)$ is fixed before $X_{i,b}$ is seen; $\sum_i w_i \le 1$ on
every path; and Assumption~\ref{c06:as-fresh} holds. Then, for any sequence of
registrations and any rule for stopping or continuing audits,
\begin{equation}\label{c06:eq-family}
\Prob\big\{\text{some coordinate with } r_i \ge \tau_i \text{ is certified}\big\}
\le \alpha .
\end{equation}
\end{proposition}

The retry problem of Section~\ref{c06:sec-retry} now has a home. A second attempt
at an audit is a new coordinate, and it needs its own share of the budget. The
analyst may retry as often as she likes, provided each retry is paid for, and the
sum of what she pays cannot exceed one. An open-ended programme can fix the share
of the $i$th coordinate in advance as $w_i = 6/(\pi^2 i^2)$, which sums to one over
all $i$, so that the threshold of the $i$th coordinate grows only as $2 \log i$.
A programme that knows part of its family can give most of the budget to that part
and hold the rest in reserve for coordinates proposed later.
In log-wealth terms the threshold for coordinate $i$ is
$A_i = \ln\bigl(\pi^2/(6\alpha)\bigr) + 2\ln i$. At $\alpha = 0.05$ the first term is
3.49, so $A_1 = 3.49$, $A_2 = 4.88$, $A_5 = 6.71$, $A_{10} = 8.10$ and $A_{20} = 9.48$.
The shares should be revisited whenever the programme changes.

How the known part should be divided is a question of cost, and the article
answers it. If attack $i$ costs $c_i$ per draw and is planned at a risk with
per-draw drift $\mu_i$, defined in Section~\ref{c06:sec-cost}, the allocation that
minimises a bound on expected total cost is $w_i \propto c_i/\mu_i$. Costly attacks,
and attacks whose risk is planned close to the tolerance, receive more of the
budget. The planning values must be fixed before any outcome is seen; a plan that
proves optimistic makes the audit longer and leaves it valid.

The article first ran the audit on Bernoulli streams of known risk, where
Assumption~\ref{c06:as-fresh} holds by construction. One coordinate had share
$w = 1/8$ of $\alpha = 0.05$, a tolerance of 0.12 and a fixed bet of 0.03, so its
threshold was 160 and its own error bound $0.00625$. At the boundary risk of 0.12,
over 2,000 runs of at most 300 attacks, it certified falsely in 0.50 per cent of
runs, below its bound. A safe stream at a risk of 0.04 was certified in 95.6 per cent
of runs after a median of 98 attacks, against 99.9 per cent for a Clopper--Pearson
test applied once after attack 300.
% source: paper/publication_versions/adaptive_audit_toy.json; paper/publication_versions/jpc_evidence.tex, Table tab:new-audit
The monitored audit gave up a little power and finished in a third of the attacks.

The more telling run was on the credit register. The committee in that experiment
declared three attacks in advance: one release with one exposure reading; six
vintages with one reading and the loan's history; and six vintages with a reading
in every live vintage and the history. It wanted the least protective floor, from
an ordered grid running from 300 down to 1, at which all three met a tolerance of
0.12. Choosing from an ordered grid by testing in a fixed order is the subject of
Chapter~\ref{ch:selection}, and requiring every declared attack to pass at a setting
is the intersection--union logic of Chapter~\ref{ch:thresholds}; the article joins
the two into a \emph{release chain}, which tests each setting only after the one
before it has passed, and spends one share of the budget on the whole chain with no
split over settings or attacks.

The true floor, estimated from a separate draw of 100,000 targets per coordinate,
was 100. There the full pooled attack succeeds at 0.116, just under the tolerance;
at 75 it succeeds at 0.133, and at 150 at 0.091.
% source: paper/publication_versions/chain_audit.json (truth.full_pooled: 0.11619, 0.13312, 0.0913; true_floor 100); jpc_evidence.tex
Over 20 independent attack streams, the release chain chose 150 in 19 and 100 in
one, and never chose an unsafe floor. Run 400 times on Bernoulli streams at the
estimated risks, it never chose an unsafe floor; with every risk raised by 0.01 it
did so in 6 runs, a rate of 0.015, within the budget of 0.05.
% source: paper/publication_versions/chain_audit.json, summary.release_chain and operating_characteristics; jpc_evidence.tex
Figure~\ref{c06:fig-designs} sets the median number of attacks each design spent
on the same 20 streams. The release chain used 6,422. Registering all 39
coordinates at the start and stopping at the first failing floor used 9,002. One
chain per attack at $\alpha/3$, which is the price of choosing the adversary after
the audit, used 13,764. The fixed-sample designs used 10,000 to 36,000.
% source: paper/publication_versions/table_chain_audit.tex; chain_audit.json summary (median_attacks 6422, 9001.5, 13764, 10000, 12000, 36000)

\begin{figure}[ht]
\centering
\begin{tikzpicture}
\begin{axis}[
  xbar, bar width=9pt, bar shift=0pt, width=0.62\textwidth, height=6.4cm,
  xlabel={Median attacks over 20 streams}, xmin=0, xmax=44000,
  symbolic y coords={Fixed sample 3000 each,Fixed sample 1000 each,Looks every 500,One chain per attack,All 39 at start,Release chain},
  ytick={Fixed sample 3000 each,Fixed sample 1000 each,Looks every 500,One chain per attack,All 39 at start,Release chain},
  yticklabel style={font=\small},
  scaled x ticks=false, xtick={0,10000,20000,30000,40000},
  xticklabel style={/pgf/number format/1000 sep={\mbox{,}},font=\small},
  nodes near coords, nodes near coords style={font=\scriptsize,/pgf/number format/1000 sep={\mbox{,}}},
  legend style={at={(0.97,0.97)},anchor=north east,draw=none,font=\small},
  area legend]
\addplot[fill=white, draw=black, postaction={pattern=north east lines}]
  coordinates {(36000,Fixed sample 3000 each) (12000,Fixed sample 1000 each) (10000,Looks every 500)};
\addplot[fill=gray!70, draw=black]
  coordinates {(13764,One chain per attack) (9002,All 39 at start) (6422,Release chain)};
\legend{Fixed-sample designs, Sequential designs}
\end{axis}
\end{tikzpicture}
% source: paper/publication_versions/table_chain_audit.tex
\caption{Attacks spent choosing a floor on the synthetic credit register, tolerance
0.12, $\alpha = 0.05$. Hatched light bars are designs that test after a planned
number of attacks; solid dark bars are sequential designs that may stop at any
attack. ``One chain per attack'' runs a separate chain for each declared attack at
$\alpha/3$; ``All 39 at start'' registers every floor and attack before the audit.
All six see the same draws, and none chose an unsafe floor.}
\label{c06:fig-designs}
\end{figure}

The chain chose 150 and not the true floor of 100 for a reason the next section
makes exact. At 100 the full pooled attack succeeds at 0.116, so close to 0.12 that
no feasible audit could certify it reliably. The cost of the sequential audit is
not paid in error. It is paid in the floor: the audit settles on a setting it can
certify in a few thousand attacks, and that setting keeps 5,553 of the 27,167
loans in the first release.
% source: paper/publication_versions/chain_audit.json, released["150"] 5553, released["1"] 27167

\section{What monitoring costs}\label{c06:sec-cost}

The length of an audit depends on how far the true risk lies below the tolerance.
For a fixed bet $q$ and a true risk $p$, each attack adds to $\log E_b$ an
increment whose mean is
\begin{equation}\label{c06:eq-drift}
\mu(p,q) = p \log\frac{q}{\tau} + (1-p)\log\frac{1-q}{1-\tau} .
\end{equation}
When $\mu > 0$ the log-wealth drifts upwards and the audit will end; when $\mu \le
0$ it may never end.

\begin{proposition}[Expected audit length]\label{c06:prop-length}
Let the outcomes be independent Bernoulli$(p)$ with $p < \tau$, let the bet be fixed
at $q \in (0,\tau)$ with $\mu(p,q) > 0$, write $A = \log\{1/(\alpha w)\}$ for the
log-threshold and $u = \log\{(1-q)/(1-\tau)\}$ for the gain from one failure, and let
$N$ be the number of attacks to certification. Then
\[
\frac{A}{\mu(p,q)} \le \E N \le \frac{A + u}{\mu(p,q)} .
\]
Over $q$, $\mu(p,q)$ is largest at $q = p$, where it equals the Kullback--Leibler
divergence between Bernoulli laws with parameters $p$ and $\tau$.
\end{proposition}

Near the tolerance that divergence is about $(\tau - p)^2/\{2\tau(1-\tau)\}$, so
halving the margin between the risk and the tolerance roughly quadruples the
audit. At a tolerance of 0.12 and an error share $\alpha w = 1/60$, an attack whose
true risk is 0.10 needs about 2,100 draws with the best fixed bet, and one at 0.06
needs about 200.
% source: paper/publication_versions/jpc_main.tex, after Proposition prop:length; checked with privacy_risk.adaptive_audit.expected_length_bound (2065 and 204)
The bound is close to the best possible for any valid test, so no cleverer
procedure escapes it. An auditor does not know $p$, and the credit-register run
used a bet that learns it, the shrunk running rate $(S_b + 1/2)/(b + 1)$ capped at
0.11. A cap set below the true risk turns the drift negative. That is what happened
at the floor of 100, where the risk of 0.116 lies above the cap: the chain could
certify that floor only by chance, as it did in one stream, and even the best fixed
bet would need tens of thousands of attacks on average.
% source: paper/publication_versions/jpc_evidence.tex, paragraph "Audit length"

So the cost of monitoring is real, and it has three parts. The first is the wider
interval of the confidence sequence, small when the risk is far from the
tolerance. The second is the larger zero-success count of
Table~\ref{c06:tab-anytime}, which buys the ability to continue after a success.
The third is the one Proposition~\ref{c06:prop-length} describes, and it is shared
with every valid procedure: a release whose risk sits just under its tolerance is
expensive to certify by any method. An institution facing that cost can spend more
attacks or choose a stricter setting. It cannot buy the certificate by looking more
often.

\section{When the audit process is the target}\label{c06:sec-process}

Every guarantee in this chapter rests on two conditions that the mathematics
cannot check: that each attack runs on a target drawn afresh, and that each bet is
placed before its outcome is seen. Both are properties of how the audit is
conducted, and an operator under pressure can break either, through carelessness,
a deadline or intent.

The first fails in quiet ways. Suppose the same target is attacked again after a
failure. Let a hidden fact about the target decide the attack, so that every repeat
gives the same result: the success probability across targets may be one half, but
after one failure on this target the chance of a later success is zero. An auditor
who treats the repeats as fresh draws can drive her evidence as high as she
pleases. The same happens if two interleaved audits share a hidden feature of their
targets, or if the adversary's information is changed after outcomes are seen.
Assumption~\ref{c06:as-fresh} fails in each case, and no bound follows.

The second fails when the operator has a choice the record does not show. The
article measured three such choices on Bernoulli streams at a risk equal to the
tolerance of 0.12, with one coordinate at the full level. An operator who can pick
among seeds, or who runs shadow audits outside the record and files only the one
that passed, turns a false-certification rate of 0.038 into 0.322 with ten tries and
0.979 with a hundred. An operator who chooses each bet after seeing the outcome
certified an attack with true risk 0.20 in every one of 2,000 runs.
% source: paper/publication_versions/table_audit_process.tex; jpc_main.tex, Section sec:security
These are the retry experiment of Section~\ref{c06:sec-retry} in another form: the
error lives in attempts the record does not hold.

The repository implements three controls that close these routes when the registry
that holds them is trusted. Targets are derived from a secret the operator commits
to at registration, combined with a later round of a public randomness beacon, so
neither the operator nor the beacon alone chooses them. Each bet is written to the
ledger before the registry runs the attack and records its outcome. The ledger is
hash-chained, so an entry cannot be changed, dropped or moved without breaking the
chain. With these in place the measured rates return to the honest ones. The
controls are a library, and a library is not a security boundary: code in the same
process can still change private state, and the hash chain detects a rewrite only
if its head is published or time-stamped outside the operator's control. An
institution that needs protection against its own operator must run the registry as
a separate service, with the attack code signed by someone other than the analyst
and every attack on the confidential register passing through it
\citep{saltzer1975}.

Some weeks later Grace, from internal audit, read the record of the Thursday
decision. She did not recompute the bound. She checked that the coordinate had been
registered before its first outcome, that every bet in the ledger preceded the
outcome it was placed on, that the chain of hashes was unbroken from the head the
Bank had published, and that the abandoned Tuesday run appeared in the ledger with
the share of the budget it had spent. Each check was a question about the process,
and each could be answered from the record without trusting the analyst.

\section{The answer}\label{c06:sec-answer}

The chapter asked what an audit must fix in advance so that its analyst may look
at will, stop when she likes and retry when she must. The answer has four parts.
The audit must fix, before the first outcome, the hypothesis it will reject, which
is that the release is unsafe; the tolerance; the attack and the law from which its
targets are drawn; and the share of the error budget it spends. It must use a
measure of evidence that is a test supermartingale under that hypothesis, either a
betting process with bets placed before outcomes or a confidence sequence. It must
treat every retry and every new attack as a new coordinate that pays for itself out
of a budget whose total is fixed. And it must keep a record, outside the analyst's
control, of every attempt, every bet and every outcome. With these four in place the
analyst may look after every attack, and the certificate means what it says. Without
them, each look can be exact and the decision can still be wrong half the time.

\begin{inpractice}[In practice: watching an audit]\sloppy
Before a simulated audit starts, write down the tolerance, the attack and its
target law, the error level and, if the programme will run more than one audit,
the share of the level this audit spends. Use an anytime-valid method if anyone
will look at results before the planned end. Record every audit attempt, including
those abandoned, and treat a retry as a new audit with its own share of the
budget. Draw targets from randomness the analyst cannot choose, and log each bet
before its outcome. Report the confidence-sequence bound as the risk figure and the
crossing of the betting process as the verdict, in separate fields. When a release
sits just under its tolerance, expect a long audit and consider a stricter setting.

\smallskip\noindent\textit{In the software.} The confidence sequence of
Proposition~\ref{c06:prop-cs} is in \texttt{privacy\_risk.sequential}
(\texttt{log\_e\_process}, \texttt{confidence\_sequence}, \texttt{anytime\_upper},
and \texttt{certification\_replicates} for Table~\ref{c06:tab-anytime});
\texttt{monitoring\_study} runs the 20-look experiment. The registered audit is in
\texttt{privacy\_risk.adaptive\_audit}: \texttt{AdaptiveAudit.register} spends a
share of the budget and refuses one that would overspend it,
\texttt{register\_release\_chain} opens a release chain, and
\texttt{AuditCoordinate.observe} advances the betting process, refusing a bet
outside $[0,\tau)$ and any outcome after certification. \texttt{zero\_success\_budget},
\texttt{expected\_length\_bound} and \texttt{drift\_weights} give the cost results.
\texttt{commit}, \texttt{target\_seed}, \texttt{CommittedCoordinate} and
\texttt{AuditLedger} are the controls of Section~\ref{c06:sec-process}; the
coordinate itself accepts outcomes from any caller, so the committed path binds only
when a registry enforces it. The retry experiment is \texttt{audit\_restarts} in
\texttt{scripts/injector\_experiment.py}, and the credit-register runs are
\texttt{scripts/adaptive\_audit\_comparison.py} and
\texttt{scripts/adaptive\_audit\_chain.py}. The update at the centre of the audit is
short:
\begin{lstlisting}
success = int(success)
self._log_e += (math.log(alternative / self.tolerance) if success and alternative
                else -math.inf if success else
                math.log1p(-alternative) - math.log1p(-self.tolerance))
self._trials += 1
self._successes += success
self._log.append((success, alternative))
self._crossed = self._log_e >= self._threshold
\end{lstlisting}
\end{inpractice}

\section{Proofs}\label{c06:sec-proofs}

\begin{proof}[Proof of Proposition~\ref{c06:prop-retry}]
For the first part, the release escapes certification only if every one of the $n$
audits fails to certify. The audits are independent and each fails with
probability $1-f$, so the probability that none certifies is $(1-f)^n$ and the
probability that some audit certifies is $1-(1-f)^n$. Since $0 < 1-f < 1$, this
tends to one as $n \to \infty$.

For the second part, imagine that the targets for all $n$ audits are drawn in
advance, whether or not the analyst will run them, and let $A_j$ be the event that
audit $j$, applied to its own draw at its prearranged level $\alpha_j$, certifies.
Because each audit is a valid level-$\alpha_j$ test of a hypothesis that is true,
$\Prob(A_j) \le \alpha_j$. The release is certified only if some audit that was
actually run certifies, and that event is contained in $A_1 \cup \dots \cup A_n$
whatever rule decided which audits to run. By the union bound,
$\Prob(\text{certified}) \le \sum_j \Prob(A_j) \le \sum_j \alpha_j \le \alpha$. No
step uses independence between the audits.
\end{proof}

\begin{proof}[Proof of Proposition~\ref{c06:prop-ville}]
Fix $c > 0$ and let $T = \inf\{b : E_b \ge c\}$, with $T = \infty$ if the set is
empty. For each $n$, the stopped process $E_{b \wedge T}$ is a nonnegative
supermartingale, so $\E[E_{n \wedge T}] \le E_0 = 1$. On the event $\{T \le n\}$ we
have $E_{n \wedge T} = E_T \ge c$, and $E_{n \wedge T} \ge 0$ elsewhere, hence
$c\,\Prob(T \le n) \le \E[E_{n \wedge T}] \le 1$. Letting $n \to \infty$ and using
continuity of probability gives $\Prob(T < \infty) \le 1/c$, and $\{T < \infty\}$
is the event $\{\sup_b E_b \ge c\}$.
\end{proof}

\begin{proof}[Proof of Proposition~\ref{c06:prop-gate}]
Fix a coordinate $i$ and work on the event $H_i = \{r_i \ge \tau_i\}$. By the
measurability conditions, the coordinate's scenario, tolerance and share are known
at its registration time $\sigma_i$, so $H_i$ and $w_i$ are determined by the
information at $\sigma_i$. Between its own draws the process $E_{i,b}$ is constant.
At draw $b$, the bet $q_{i,b}$ is known beforehand and, by
Assumption~\ref{c06:as-fresh}, the conditional success probability $p$ satisfies
$p \ge \tau_i$, so by \eqref{c06:eq-factor} the conditional expectation of the
update factor is at most one. Hence, on $H_i$, the process indexed by the global
sequence of recorded events and started at $\sigma_i$ is a nonnegative
supermartingale starting at one. Applying Proposition~\ref{c06:prop-ville}
conditionally on the information at $\sigma_i$, with $c = (\alpha w_i)^{-1}$, gives
$\Prob(\text{$i$ certified} \mid \mathcal{F}_{\sigma_i})\,\ind{H_i} \le \alpha w_i
\ind{H_i}$; when $w_i = 0$ the threshold is infinite and the left side is zero. Taking
expectations, $\Prob(\text{$i$ certified},\ H_i) \le \alpha\,\E[w_i]$. Summing over
$i$, using the union bound and monotone convergence for countably many
coordinates, and then $\sum_i w_i \le 1$ on every path, gives
$\Prob(\exists i : \text{$i$ certified},\ H_i) \le \alpha\,\E[\sum_i w_i] \le
\alpha$, which is \eqref{c06:eq-family}.
\end{proof}

\begin{proof}[Proof of Proposition~\ref{c06:prop-cs}]
For fixed $u \in (0,1)$, the likelihood ratio after $B$ attacks of a success rate $u$
against the rate $r_0$ is $L_B(u) = (u/r_0)^{S_B}\{(1-u)/(1-r_0)\}^{B-S_B}$. Given
the past, the next factor has expectation $r_0 (u/r_0) + (1-r_0)(1-u)/(1-r_0) = 1$
under conditionally Bernoulli$(r_0)$ outcomes, so $L_B(u)$ is a nonnegative
martingale with $L_0(u) = 1$. Let $\pi$ be the beta$(a,b)$ density. By Tonelli's
theorem, $\int L_B(u)\pi(u)\,du$ is again a nonnegative martingale starting at one,
and evaluating the integral,
\[
\int_0^1 \frac{u^{S_B}(1-u)^{B-S_B}}{r_0^{S_B}(1-r_0)^{B-S_B}}
\frac{u^{a-1}(1-u)^{b-1}}{\mathrm{B}(a,b)}\,du
= \frac{\mathrm{B}(a+S_B,\, b+B-S_B)}{\mathrm{B}(a,b)\,r_0^{S_B}(1-r_0)^{B-S_B}}
= E_B(r_0) .
\]
By Proposition~\ref{c06:prop-ville} with $c = 1/\alpha$, the probability that
$E_B(r_0) \ge 1/\alpha$ for some $B$ is at most $\alpha$. On the complementary event
$r_0 \in \mathcal{C}_B$ for every $B$, which proves the claim.
\end{proof}

\begin{proof}[Proof of Proposition~\ref{c06:prop-length}]
Write $Y_b$ for the increment of $\log E_b$ at attack $b$. The $Y_b$ are independent
with mean $\mu = \mu(p,q) > 0$, and each is at most $u$, the value taken on a
failure. Fix a cap $B$ and let $T = \min(N, B)$, a bounded stopping time. Wald's
identity gives $\E[\log E_T] = \mu\,\E T$. If $T = N$, then $\log E_{N-1} < A$ and
$\log E_N < A + u$; if $T = B < N$, then $\log E_B < A$. In both cases
$\log E_T < A + u$, so $\mu\,\E T < A + u$, that is $\E\min(N,B) \le (A+u)/\mu$.
Letting $B \to \infty$ and using monotone convergence gives $\E N \le (A+u)/\mu$; in
particular $\E N < \infty$, so Wald's identity applies to $N$ itself. Since
$\log E_N \ge A$ at certification, $\mu\,\E N = \E[\log E_N] \ge A$, which is the
lower bound. Finally, $\partial \mu / \partial q = p/q - (1-p)/(1-q)$, which is
positive for $q < p$ and negative for $q > p$, so $\mu$ is maximised at $q = p$,
where $\mu(p,p) = p\log(p/\tau) + (1-p)\log\{(1-p)/(1-\tau)\}$, the
Kullback--Leibler divergence of Bernoulli$(\tau)$ from Bernoulli$(p)$.
\end{proof}
